\documentclass[11pt]{article}
\usepackage[UTF8]{ctex}
\usepackage{booktabs,amsmath,amssymb,graphicx}
\usepackage[margin=1in]{geometry}
\usepackage[hidelinks]{hyperref}
\title{TLI：面向长上下文稀疏注意力的免训练轻量索引\\
\large 块级上界粗筛与远/近分区预算}
\date{2026-09-30}
\begin{document}
\maketitle

% =====================================================================
% Abstract：模仿 MoBA 四句式（stakes+障碍 → 两类现有方法批判 →
% we propose+命名+key advantage → 部署声明）。
% =====================================================================
\begin{abstract}
扩展大语言模型的有效上下文长度需要突破 KV 缓存带宽的制约：在 131K 上下文的 8B 模型上，decode 每步的全量 KV 读取远超计算量，成为吞吐上限。现有方案面临一个困难权衡：training-free 的稀疏索引依赖下界近似，在长上下文检索任务上质量崩塌（我们实测 Quest 在 RULER 16K 的 multikey\_3 上仅 21 分，FullKV 为 100）；训练类索引（DSA）引入可学习参数与上千步 warm-up 成本，且选择质量依赖训练分布。本文提出 TLI，一个两级 training-free 稀疏注意力索引：第一级用 4bit 量化子空间上的块级 minmax \textbf{上界}粗筛保证子空间口径下无漏选，第二级用 far/near 分区的 token 级精筛保护远端关键 token 不被近端高分挤出。在 Qwen3-8B 的 LongBench 13 任务上，TLI 取得 50.54，与 FullKV（50.36）在噪声级差距内持平；多跳 far-heavy 任务 musique 达到 34.76，为全部方法（含 FullKV）最高。RULER 上 87.93 逼近 FullKV 的 88.71，分区相对单池消融臂的增益在 16K 才转正（4K/8K 为 $-$0.14/$-$0.25，16K 为 +0.91），增益集中于多干扰检索与词表任务。速度侧，TLI 选择链相对 dense 打分取得 1.46$\times$@32K $\rightarrow$ 5.09$\times$@128K 的 kernel 级加速，64K 单请求 e2e 1.285$\times$；全部精度代价与负结果在同机同协议下如实报告。TLI 已在 sglang 生产形态（paged 寻址、增量索引、CUDA graph）下集成。
\end{abstract}

% =====================================================================
% 1 Introduction：模仿 MoBA 六段式。
% P1 重要性+应用锚点 → P2 技术挑战+方向 → P3 三类现有方法批判
% → P4 设问 → P5 we introduce+机制 → P6 贡献列表。
% =====================================================================
\section{引言}

% P1（stakes）：长上下文是当代 LLM 的核心能力，KV 带宽是结构性瓶颈。
处理长上下文的能力正在成为大语言模型的核心竞争力。DeepSeek-R1~\cite{dsr1} 的长链推理、Kimi 与 Claude 的百万级上下文产品、以及长文档问答与代码仓库理解，都要求模型在一次请求中消费数万至数十万 token。然而上下文的每次延长都直接推高 KV 缓存的体量：131K 上下文配 8B 模型的 KV 读取已远超每步计算量，decode 阶段的全量 KV 访问构成带宽瓶颈，prefill 阶段同样受全注意力带宽制约。推理系统侧的进步（PagedAttention~\cite{vllm}、FlashAttention~\cite{fa}）优化了显存组织与计算核，但没有减少需要读取的数据量本身。

% P2（挑战+方向）：attention 稀疏性是减数据量的出路；三层分布是本文的观察起点。
减少 KV 读取量的出路在于注意力的稀疏性：softmax 归一化使少数 token 承载绝大部分 attention mass。我们在 Qwen3-8B 真实数据上的 trace 分析给出这个分布的具体形态——dense top-1024 选择的 mass 由三层构成：sink（0.37--0.71）、近端滑窗（0.16--0.29）、远端检索（0.02--0.29），且远端层轮廓在同任务内跨层相关高达 0.93--1.00。这个三层结构意味着：稀疏索引要保住精度，必须同时喂饱三类性质完全不同的 attention，而它们的预算需求互相冲突。

% P3（三类批判）：每类给机制+定量缺陷。
现有方案在三类路径上各自受限。第一类，training-free 稀疏索引以 Quest~\cite{quest} 为代表，用 page 内采样的最小值构造分数\textbf{下界}近似——下界会低估真实高分页，干扰 key 随长度超线性增多时漏选加剧，我们实测其在 RULER 16K multikey\_3 上崩至 21 分（FullKV 100）；KV 压缩方法（SnapKV/H2O/PyramidKV~\cite{snapkv,h2o,pyramidkv}）在 prefill 时一次性丢弃 KV，无法响应 decode 侧的查询分布。第二类，训练类索引以 DSA~\cite{dsa}（DeepSeek V3.2）为代表，用可学习投影与 ReLU 打分换选择能力，需要 1000 步 warm-up 与 2.1B token 的训练预算，且选择质量依赖训练分布。第三类，两级结构（HISA~\cite{hisa}，COLM 2026）验证了块级粗筛加 token 级精筛的形态有效，但其中两个正交设计维度——分数表示取上界还是下界、预算组织取分区还是全池——对质量与速度的影响尚无人系统刻画。

% P4（设问，MoBA 式）：critical research question。
由此产生一个关键问题：能否设计一个 training-free、零 warm-up 的稀疏注意力索引，在不引入任何可学习参数的前提下，同时保住长上下文检索质量并兑现带宽收益？

% P5（we introduce + 机制完整段）。
本文提出 TLI（Training-free Lightweight Indexer）回答这个问题。TLI 是一个两级级联的索引：第一级把 KV 缓存组织为 64-token 块，在 position-stable 低频子空间（$d'{=}32$）上维护每块的 min/max 上界并以 4bit 量化存储，query 打分取块上界的乐观侧——上界保证无漏选，任何真实高分 token 所在块必然入选；第二级把 token 预算划分为 far/near 两个池，近端滑窗强制覆盖局部性，远端池对 L1 入选块做 token 级精筛，独立配额使远端关键 token 不被近端高分挤出。两级之间以感知型 per-request gate 收尾：prefill 末段的一次 softmax 统计识别 far mass 近零的层，decode 侧安全跳过其远端检索。整个选择链 kernel 化为 fused L1 块打分、L2 级联 dual topk 与统一稀疏 attention kernel，O(n) 增量索引使块界维护与序列长度无关。

% P6（贡献列表，claim-first）。
本文的贡献如下：
\begin{enumerate}
\item \textbf{training-free 降维自由度的系统刻画}。位置稳定低频子空间在索引质量上无损失（子空间选择对全维 top-1024 的 mass 覆盖为 0.729，全维打分自身 0.732，随机子空间 0.32）；tail32 的机理经三重判决确证为\textbf{16 个完整最低频 RoPE 旋转对}（far 全链 mass recall 口径：频率单调——低频对 0.811 / 中频对 0.360 / 高频对 0.004；配对完整性——错位配对崩至 0.567；16 对即饱和——8 对仅 0.473）。选取方法学四象限给出可行域边界：逐维谱（0.759）与训练投影均败于配对纠缠；query 自适应 pair 选取（0.865）不可索引化；per-layer 静态 pair（0.849）兼容索引——tail32 是零信号成本最优解。无监督 PCA 可将细筛再压至 $d{=}16$ 近无损（0.766 vs 0.804），且跨 kv-head 共享 SVD 基在 8/8 样本上追平 tail32；而训练投影在 trace 重放中成立（92--94\%）、e2e 崩坏——投影基对 decode 查询分布漂移脆弱，旋转对频率先验与查询分布无关恒稳。降维自由度来自位置先验，而非学习或自适应。
\item \textbf{far/near 分区预算}。选择索引建在 kv-head 级：同组 query head 共享同一份 kv-head 级选择，无需聚合。GQA 下远端 attention mass 在 kv-head 间极不均（单 kv-head 可独占 0.804），全池 top-k 使远端关键 token 被近端高分挤出。分区独立配额在双基准 e2e 上成立（LongBench +0.25 / RULER +0.17 vs 单池），多跳 far-heavy 任务增益最纯粹（musique 34.76 vs 同预算单池臂 29.74，$+5.02$）；聚类代表方案在严格 token 预算下无一致优势，降级为消融负结果。
\item \textbf{感知型 per-request gate}。静态层掩码跨任务不泛化（多跳任务掉 4.8--5.9 分）；prefill 末段动态信号与 decode 侧 far 质量相关性 0.924，真实触发区间（$\tau{=}0.005$）下，开启 gate 相对关闭 gate 的差分：musique 微升 $+0.23$、qasper 微降 $-0.41$，far-heavy 任务均不崩。gate 的定位是安全省算力的保守开关，而非精度增益点。
\item \textbf{免调参的轻量配置选择}。分区参数 $\alpha/\beta/\gamma$ 的质量面在任务与层两个粒度上都高度平坦：逐任务最优选择的 oracle 上限仅 $+0.08$（13 任务）、逐层最优仅 $+0.0008$，prefill 端一次 far 统计驱动的逐样本在线选择净增益为 0。因此部署只需一组默认配置，far-heavy 多跳任务可选一次 $\beta$ 上调；参数不需精调本身构成部署优势（对比 DSA 的千步 warm-up）。
\end{enumerate}

% =====================================================================
% 2 Method：模仿 MoBA §2 结构——Preliminaries 形式化 → 架构+核心公式
% → 因果性两设计 → 设计选择讨论（子空间/分区预算）→ 与 sink/SWA 关系
% → Implementation。
% =====================================================================
\section{方法}

\subsection{预备：标准注意力与稀疏选择问题}

% MoBA 2.1 式：单 head 形式化，简洁。
先回顾标准注意力。单个查询 token $q\in\mathbb{R}^{1\times d}$ 访问 $N$ 个 KV token（$K,V\in\mathbb{R}^{N\times d}$）：
\begin{equation}
\mathrm{Attn}(q, K, V) = \mathrm{Softmax}\!\left(qK^\top\right)V.
\end{equation}
稀疏索引的问题是：为 $q$ 选出大小为 $K_2$ 的位置集合 $I\subseteq[N]$，使 $\mathrm{Attn}(q, K[I], V[I])$ 逼近全量结果。本文关注单 head 情形以保持清晰；GQA 下逐 kv-head 打分与远端 mass 的跨 head 不均性，在 §~\ref{sec:l2} 的分区机理中讨论。与训练类索引不同，training-free 索引不允许任何参数——选择质量完全取决于用什么信号打分、以什么组织方式分配预算。

\subsection{TLI 架构}\label{sec:l2}

\begin{figure}[h]\centering
\includegraphics[width=0.85\linewidth]{figures/fig7_tli_architecture.pdf}
\caption{TLI 两级索引架构：L1 块级 minmax 上界粗筛（子空间 $d'{=}32$ + 4bit 量化）$\rightarrow$ L2 far/near 分区 token 级精筛 $\rightarrow$ 统一稀疏 attention fused kernel。}
\label{fig:arch}
\end{figure}

% MoBA 2.2 式：Different from standard attention where..., X enables...
与标准注意力每个查询访问全部上下文不同，TLI 为每个查询 head 选出 $K_2{=}1024$ 个 KV token（图~\ref{fig:arch}）：
\begin{equation}
\mathrm{TLI}(q, K, V) = \mathrm{Attn}\!\left(q,\ K[I_{\mathrm{L1}\rightarrow\mathrm{L2}}],\ V[I_{\mathrm{L1}\rightarrow\mathrm{L2}}]\right),
\end{equation}
其中 $I_{\mathrm{L1}\rightarrow\mathrm{L2}}$ 由两级级联产生。TLI 的关键设计不在级联形态本身（HISA 已验证），而在两个正交维度上的选择：\textbf{分数表示取上界}与\textbf{预算组织取分区}。

\textbf{L1：块级上界粗筛。}把上下文按 $B{=}64$ 切块，块 $I_i$ 内 token 记为 $k_j$。对每块每 kv-head，在 position-stable 低频子空间（$d'{=}32$）上维护逐维 min/max（$k^{\min}, k^{\max}\in\mathbb{R}^{d'}$）并以 4bit 量化存储（128 维压至 40B/token-head）。查询对块的打分取上界的乐观侧：逐维按 $q_d$ 的符号取 min/max 端值求和，构成该块对块内任意 token 点积的上界
\begin{equation}
s_{\mathrm{blk}}(q, i) \;=\; \sum_{d=1}^{d'} \max\!\big(q_d\, k^{\min}_{i,d},\; q_d\, k^{\max}_{i,d}\big)\;\geq\; \max_{j\in I_i}\; \langle q, k_j \rangle.
\end{equation}
这个上界在\textbf{子空间打分口径下}保证无漏选：真实高分 token 所在块的子空间分数必然不低于该 token 的子空间真实分数，必然进入 top-$K_1$（子空间相对全维点积是有损的，损失由 §2.3 的选取方法学单独刻画）。这与 Quest 的 page-min 下界近似形成对照——下界会低估真实高分页，漏选随干扰项增多而加剧。取 top-$K_1{=}128$ 块进入 L2。

\textbf{L2：far/near 分区精筛。}sink 与滑窗（最近 $sw_{lo}$ 个 token，局部性由滑窗语义天然保证，无需打分）构成强制保留区、不占分区预算；剩余 mid 预算划分为两个池。far 池对 L1 入选块做 token 级 4bit 精筛，取 far 配额 $F$（128--256 饱和）；near 池对近端区做同口径精筛。分区防挤出的机理：GQA 下远端 mass 在 kv-head 间极不均（trace 实测单 head 独占 0.804 而同组其余 head 近零），全池 top-k 时 far 候选与近端高分 token 在同一池竞争，far-heavy 层的远端关键 token 会被系统性挤出。分区给 far 独立配额，近端扣减带 near\_floor 保护防边界回退。

\textbf{预算的形式定义。}设查询前缀长度 $S=t+1$、mid 区长度 $L_{\mathrm{mid}}=S-|I_{\mathrm{sink}}|-sw_{lo}$、近端区长度 $\ell_{\mathrm{near}}=\lceil\alpha L_{\mathrm{mid}}\rceil$。两池与最终选择集为
\begin{align}
I_{\mathrm{far}} &= \operatorname*{TopK}_{F}\Big\{\langle q,k_j\rangle:\ j\in\bigcup\nolimits_{i\in I_{\mathrm{L1}}} I_i\cap\big[\,|I_{\mathrm{sink}}|,\ S-sw_{lo}-\ell_{\mathrm{near}}\,\big]\Big\},\\
I_{\mathrm{near}} &= \operatorname*{TopK}_{K_2^{\mathrm{mid}}-F}\Big\{\langle q,k_j\rangle:\ j\in\big[\,S-sw_{lo}-\ell_{\mathrm{near}},\ S-sw_{lo}\,\big]\Big\},\\
I_{\mathrm{L1}\rightarrow\mathrm{L2}} &= I_{\mathrm{sink}}\ \cup\ \big[\,S-sw_{lo},\ t\,\big]\ \cup\ I_{\mathrm{near}}\ \cup\ I_{\mathrm{far}},
\end{align}
其中 mid 预算 $K_2^{\mathrm{mid}}=K_2-|I_{\mathrm{sink}}|-sw_{lo}$，far 配额 $F=\max\big(64,\ K_2^{\mathrm{mid}}-\lfloor\gamma\,\lceil\ell_{\mathrm{near}}/B\rceil B\rfloor\big)$：$\gamma$ 经近端块数折算近端细筛配额，far 取剩余预算并保底 64；$\beta$ 作用于 L1 的近端页预算。

\textbf{维持因果性的两个设计。}自回归语言模型要求查询不能访问未来 token。TLI 通过两个具体设计维持因果性：其一，\textbf{池边界即因果边界}——far 池的上界位置 $\leq t+1-\ell_{\mathrm{near}}-sw_{lo}$、near 池 $< sw_{lo}$，两池的并集天然落在因果范围内，无需显式因果 mask（L1 的垃圾块也天然落池外）；其二，\textbf{早期行均匀覆盖}——首个 prefill chunk 内，行位置 $t_r$ 小于 $K_2$ 的早期行候选不足，若放任 topk 返回的哨兵/未来位置混入，非因果泄漏会逐层传播（我们在 30B 模型上实测到生成重复循环崩坏，根因正是此处；8B 对同 bug 耐受未崩，如实报告）。早期行改用均匀重复 grid，等价 dense。

\textbf{与 sink 和滑窗注意力的关系。}sink（StreamingLLM~\cite{streamingllm}）与滑窗注意力（Longformer/BigBird~\cite{longformer,bigbird}）是两种成熟的静态稀疏结构，本文将它们作为已知组件沿用：sink token 承载 0.37--0.71 的 attention mass，滑窗承载近端局部性——两者共同构成 TLI 的强制保留区，不占分区预算。本文的创新不在保留区本身，而在\textbf{剩余预算在 far 检索与 near 精筛之间的分配}。作为旁证，我们在 baseline 适配中实测：SnapKV/H2O/PyramidKV 无显式 sink 保送时在 Qwen3-8B 上首 token 即 EOS（末层 sink 权重 mean 0.592 而浅层选择信号不指向 sink，逐层恶性畸变）——任何不显式处理 sink 的稀疏方案在该模型族上都会失真。

\subsection{设计选择讨论}

% MoBA 式：Next, we discuss some additional key design choices...
接下来讨论 TLI 的三个关键设计选择。

\textbf{子空间选择。}L1 为什么取低频尾维 32 维。索引信号必须在降维后保住对全维分数排序的区分力。Qwen3 的 rotate\_half RoPE 使低频维上的点积与全维点积排序高度一致（子空间选择对全维 top-1024 的 mass 覆盖 0.729，全维打分自身 0.732），随机子空间崩至 0.32、高频维 0.137。机理经三重判决确证（图~\ref{fig:dim}）：Qwen3 的 128 维全部参与 RoPE，rotate\_half 布局下频率 $j$ 施加在维度对 $(j, j{+}64)$ 上，tail32 = 维度 48..63 与 112..127 = \textbf{16 个完整最低频旋转对}。点积按频率分解为 $\sum_j [(q_j k_j + q_{j+64} k_{j+64})\cos(\Delta\,\omega_j) + (q_j k_{j+64} - q_{j+64} k_j)\sin(\Delta\,\omega_j)]$，远端大 $\Delta$ 上低频项平滑高相关、高频项震荡相消——低频对承载长程检索信号。三个判决支持这个结构解释：其一，频率单调（低频 16 对 0.811 / 中频 0.360 / 高频 0.004）；其二，配对完整性（前半取低频 16 维、后半换成中频对来源的错位配对崩至 0.567，段位几乎重合仅元素来源错位）；其三，对数饱和（16 对即 0.811 近饱和，8 对仅 0.473）。32B 上方向复现（完整对 0.829 > 错位 0.645 > 单半 0.54），饱和宽度后移至 24 对（0.904）。

选取方法学的完整刻画（表~\ref{tab:select}）：静态（索引兼容）与自适应（逐 query）两个可行域交叉出四个象限。逐维谱 top-32 仅 0.759——旋转对两元素在点积中纠缠，逐维重要性原理上无法恢复配对结构。SparQ~\cite{sparq} 式逐 query $|q|$ 幅值 top-32 选取达 0.848，pair 感知自适应（按 $|q_j|+|q_{j+64}|$ 取完整对）达 0.865——自适应选取确有小增益，但其子空间逐 query 变化，与 query 到达前建好的增量 minmax 块索引结构性不兼容（SparQ 为此付出逐 query 维度随机读取代价，不建索引）。中间点是 per-layer 静态化：prefill 末段 $q$ 统计的 pair 幅值 top-16 完整对，层固定、可索引化，0.849——保住自适应增益的七成，比 tail32 高 3.8pt，代价是每层一次 $q$ 统计。\textbf{tail32 的准确定位是零信号成本可行域的最优解}，而非全局最优；pair 静态选取是可索引化的升级路径。训练投影可以做得更准（trace 重放 92--94\%），但 e2e 崩坏：投影基对 decode 查询分布漂移脆弱，而旋转对频率先验与查询分布无关恒稳。这条 trace 成立而 e2e 崩坏的口径鸿沟，是 training-free 与训练流派之间分界线的实证刻画。

\begin{figure}[h]\centering
\includegraphics[width=0.95\linewidth]{figures/fig2_dim_reduction_story.pdf}
\caption{降维自由度来自位置先验：(a) tail32 = 16 个完整最低频旋转对（频率单调、错位配对崩 0.567、16 对饱和）；(b) 共享 SVD 基 $d{=}16$ 在 8/8 样本上追平 tail32；(c) 线性优于非线性、共享优于 per-head。}
\label{fig:dim}
\end{figure}

\begin{table}[h]\centering
\caption{子空间选取方法学四象限（far 全链 mass recall，8 样本均值）。}
\label{tab:select}
\begin{tabular}{llcc}
\toprule
选取方法 & 信号 & recall & 索引兼容 \\
\midrule
逐维谱 top-32 & 数据驱动（静态） & 0.759 & 兼容 \\
tail32 & 频率先验（静态） & 0.811 & 兼容 \\
SparQ 式 $|q|$ top-32 & query 自适应 & 0.848 & 不兼容 \\
pair 幅值 top-16 对 & query 自适应 & 0.865 & 不兼容 \\
per-layer 静态 pair & prefill $q$ 统计 & 0.849 & 兼容 \\
\bottomrule
\end{tabular}
\end{table}

\textbf{分区预算参数。}预算空间由三个参数刻画：$\alpha{=}0.125$（near 区比例）、$\beta{=}0.375$（near 页预算）、$\gamma{=}0.125$（near token 预算），far/near 方法取 (minmax, avg)。132 个配置点的网格加上五个方法组合臂的 e2e 给出三条规律：$\alpha$ 单调至边界（越省 far 越好，受 near\_floor 制约）；$\beta$ 在 0.25--0.75 上平坦（多跳任务 0.375 占优）；$\gamma$ 必须随 $K_2$ 等比缩放（$\gamma{=}0.5$ 在 $K_2{=}1024$ 下 far 仅保底 64，F1 崩至 13.51 的教训）。最重要的发现是平坦性本身：per-task 完美选 $\beta$ 的 oracle 上限仅 +0.08，per-layer oracle 仅 +0.0008——\textbf{分区预算对配置误差不敏感，不需要精调}。这构成部署简单性论据：对照 DSA 的 1000 步 warm-up，TLI 一组固定配置覆盖 12/13 任务。

\textbf{gate 的阈值。}动态 gate 用 prefill 末段 far 统计识别 far mass 近零的层。阈值 $\tau$ 的真判决（双任务梯度）：$\tau{=}0.005$ 下 musique +0.23 / qasper $-$0.41；$\tau{=}0.015$ 下 +0.27 / $-$0.67——梯度单调，真触发区间微损不崩（对照静态版掉 4.8--5.9）。gate 因此定位为安全省算力的保守开关。

\subsection{实现}

% MoBA 2.3 式：We provide a high-performance implementation..., consisting of N major steps.
TLI 在 sglang 生产形态下实现，选择链由四个 kernel 组件构成：
\begin{itemize}
\item \textbf{fused L1 块打分}：单个 Triton kernel 内完成 gather + 4bit 反量化（共享 scale 表）+ GEMV 打分 + 块级规约，32K 形态 3.6$\times$；
\item \textbf{L2 级联 dual kernel}：单 launch 完成 far/near 双池打分与分区 topk，池边界即因果边界（替代 eager 的 fine 矩阵因果 mask）；
\item \textbf{统一稀疏 attention fused kernel}：gather + online softmax + bf16 直写，kernel 级 11--16$\times$；
\item \textbf{O(n) 精确增量索引}：每 token 增量更新块界与全量重建逐位一致，0.128ms 与序列长度无关。
\end{itemize}
系统侧三项工程对接生产：paged 寻址（选择输出逻辑位置经 req\_to\_token 间接寻址 KV pool，零拷贝）；per-layer 共享 index pool 预分配（decode 批量 select 的 launch 数与 batch size 无关）；CUDA graph 捕获区零 host 同步。设计约束来自 H20 形态：Tensor Core 吞吐仅 H100 的 15\% 而 HBM 带宽同级——打分类 kernel 的第一性约束是带宽而非算力，TC 化与 TMA 转置布局实测均否决（算术强度指标显示纯 gather 路径已带宽饱和），kernel 主线是消除中间物化与合并 launch。

% =====================================================================
% 3 Experiments：模仿 MoBA §3 句式（目标句 → 实验 → As shown in...
% 数字 → 机理 → This suggests that...）。
% =====================================================================
\section{实验}

\subsection{端到端质量}

% MoBA 3.1 式：In this section, we conduct... to validate...
本节在双基准上验证 TLI 的质量：RULER~\cite{ruler}（3 长度 $\times$ 11 任务 $\times$ n=100，NVIDIA 官方预生成）考察长上下文检索，LongBench~\cite{longbench} 13 子集考察真实任务混合口径。对比方法中，TIA 是与本文同一管线下的第一代 training-free 块上界索引器（块 min/max 上界 + 4bit 子空间量化，无分区预算与感知 gate），作为 TLI 最直接的结构消融对照；Quest 为下界近似的代表；KV 压缩三 baseline 走统一 monkeypatch 管线。全部实验用 Qwen3-8B~\cite{qwen3} 真实权重；TLI 用终局配置（sink/swa 强制保留不计入预算的严格口径）。

\begin{table}[h]\centering
\caption{RULER 四方法终表（string\_match\_all $\times$100；TLI 为分区配置：(minmax, avg) 打分、$\alpha{=}0.125$、$\beta{=}0.25$、$\gamma{=}0.125$、$K_2{=}1024$）。}
\begin{tabular}{lllll}
\toprule
 & FullKV & TIA@1024 & TLI@1024 & Quest@1024 \\
\midrule
4K & 91.59 & \textbf{91.63} & 91.49 & 87.91 \\
8K & \textbf{89.14} & 88.90 & 88.62 & 81.76 \\
16K & \textbf{85.41} & 84.53 & 83.68 & 70.18 \\
总 AVG & \textbf{88.71} & 88.35 & 87.93 & 79.95 \\
\bottomrule
\end{tabular}
\end{table}

如表所示，三个观察成立。其一，TLI 对 Quest 的领先随长度超线性扩大（+3.58 $\rightarrow$ +6.86 $\rightarrow$ +13.50）：16K 档 Quest 在多干扰检索任务上崩塌，TLI 同任务保持量级正确——根因是下界近似漏选随干扰项增多加剧，上界粗筛无此劣化。其二，分区相对单池消融臂的增益在 16K 才转正（4K/8K 为 $-$0.14/$-$0.25，16K 为 +0.91）：超长上下文是分区 far 检索的收益区，16K 的分区增益集中于 multikey\_3（+5.0）与 cwe（+5.4），均为 vs 同配置单池臂的差分。其三，绝对值解读须以同长度 FullKV 为基线——聚合列举类任务 FullKV 自身随长度退化（multivalue 89.5 $\rightarrow$ 69.25 $\rightarrow$ 44.25 是模型列举能力上限），稀疏方法在同长度上的差距才是稀疏效应。

LongBench 13 子集上（表~\ref{tab:longbench}），TLI \textbf{50.54}——与 FullKV 50.36 持平（+0.18，噪声级边缘，Limitations 一节如实标注），vs Quest +2.82。增益的结构清晰：多跳 far-heavy 任务 musique 34.76 为全部方法最高（vs FullKV +2.62、vs 同预算单池臂 +5.02），是分区 far 检索的直接证据；其余 12 任务在 $0\sim-0.44$ 噪声级互有胜负。KV 压缩三 baseline 的全量崩坏坐实 sink 论断：SnapKV 27.70 / H2O 更崩（musique 0.63、hotpotqa 2.68）——无显式 sink 保送时 Qwen3-8B 首 token 即 EOS（§2.2 末段机理），三者中仅代码类任务（lcc 67.19 / 54.06）部分幸存。预算配比的任务形态依赖性如实并报：LongBench 多跳任务 $\beta{=}0.375$ 占优（50.54 行），RULER 合成 needle $\beta{=}0.25$ 占优（87.93 行）——两行并报，使用者按任务形态选择。逐任务族分解给出 $\beta$ 让渡的边界判据：真实任务全族成立（最 far-heavy 的 musique 反而增益最大 +1.94），合成 far-dense 族失效（multikey\_3 $-$2.00：far 候选数量爆炸使预算粒度主导）——far-dense 检索是稀疏索引器全族难区（Quest 同任务同构崩塌），非 $\beta$ 配比独有问题。

\begin{table}[h]\centering
\caption{LongBench 13 子集终表（各任务官方口径平均；KV 压缩三 baseline = 统一 monkeypatch 管线 + 1024 预算，无 sink 保送适配）。}
\label{tab:longbench}
\begin{tabular}{lccccccc}
\toprule
 & FullKV & Quest & TIA & TLI & SnapKV & H2O & PyramidKV \\
\midrule
总分 & 50.36 & 47.72 & 50.06 & \textbf{50.54} & 27.70 & 14.87 & 28.82 \\
musique & 32.14 & 27.09 & 32.28 & \textbf{34.76} & 8.61 & 0.63 & 8.99 \\
\bottomrule
\end{tabular}
\end{table}

\subsection{消融}

% MoBA 式：目标句 → 数字 → 机理 → 结论。
\textbf{子空间（A）。}低频 $d'{=}32$ 对全维 top-1024 的 mass 覆盖 0.729，全维打分自身 0.732（随机 0.32、高频 0.137 崩溃）——位置稳定子空间存在。机理三重判决：tail32 = 16 个完整最低频旋转对（频率单调 0.811/0.360/0.004、错位配对崩 0.567、16 对饱和）；选取方法学四象限（表~\ref{tab:select}）刻画可行域——逐维谱 0.759 与训练投影均败于配对纠缠，query 自适应 pair 选取 0.865 是 oracle 上界但不可索引化，per-layer 静态 pair 0.849 兼容索引，tail32 是零信号成本最优。降维三段判决：无监督 PCA $d{=}16$ 近无损（0.766 vs 0.804）是唯一推荐；训练投影 trace 成立 e2e 崩（口径鸿沟）；非线性（UMAP/t-SNE）质量无增益且 fit 成本 80--167s/层-样本 vs 增量索引 0.128ms/token，四个数量级不可行。跨 kv-head 共享 SVD 基 $d{=}16$ 在 8/8 样本追平 tail32 且优于 per-head 基——GQA head 冗余使 MLA 式单投影缓存的 training-free 代理可行。

\textbf{分区。}far 区 token 级精筛接近 oracle（far-heavy 层 0.999）；聚类代表在严格 token 预算下无一致优势（块 scatter-amax 0.09--0.39 最差）且聚类中心非上界有漏选理论空洞——降级为负结果。分区 vs 单池 e2e 双基准成立（LongBench +0.25 / RULER +0.17），与 trace 重放口径（宽预算下单池微胜）的分歧坐实：\textbf{分区收益的本质是紧预算下的预算分配，而非候选质量}。

\textbf{层跳过（D'）。}离线平均轮廓可跳 13/36 层（precision 0.92--1.00）但全局静态掩码跨任务不泛化（多跳掉 4.8--5.9）；prefill 动态信号 corr 0.924 防住反向错误，真触发区间微损不崩——gate 是安全开关。

\textbf{自适应 $\alpha/\beta$ 的否定。}per-request 在线分派 $\alpha/\beta$ 判 No-Go：oracle 上限 LongBench +0.08 / RULER +0.00，且 +0.08 全部可由噪声正向选择偏差解释（非 musique 12 任务 $|\Delta|$ 中位 0.17）。信号零成本可得，但参数面平坦使自适应无杠杆。同一信号的最佳用途是 gate（安全价值）而非 $\beta$ 分派（无精度价值）。粒度细化到层同样无杠杆（per-layer oracle 增益均值 +0.0008）——任务与层两粒度双重否定，转译为部署简单性论据。

\textbf{离线重放与端到端的口径鸿沟。}我们对五种近/远端打分组合做了离线 trace 重放与端到端评测的双重排名：重放质量排名第 4 的组合（近端取均值聚合的配置）在端到端评测中取得全场最高分 50.54——两种口径的排名系统性反转。这说明离线代理指标不能替代端到端评测，kernel microbench 与 e2e 必须都测都报告。

\subsection{Kernel 级 Microbench}

选择链相对 dense 全维打分的加速比：32K 1.46$\times$、64K 2.77$\times$、128K 5.09$\times$（合成数据已标注，trace 重放口径）。同机三方对比（官方 kernel 原样接入统一 harness，131K cudaEvent）：Quest 0.107ms / DSA 0.503ms / TLI 0.787ms。延迟排名 Quest $<$ DSA $<$ TLI，如实报告——TLI 当前的优势在算法结构侧：每 token 索引 MAC $\approx$258（DSA 的 1/32）、索引存储 3$\times$ 低于 Quest、质量 vs Quest LongBench +2.82 / RULER +7.98。131K 下三家都远离 HBM bound，排名反映实现成熟度而非算法上界。

\subsection{端到端速度}

e2e 终值（Qwen3-30B-A3B，H20$\times$2）：64K 单请求 1.285$\times$（16.39s vs dense 21.06s）——稀疏理论流量收益首次在 e2e 净兑现；32K 档 0.77$\times$（短上下文固定开销未能覆盖）；TP2 bs16 64K 档 0.92$\times$，差 8\% 如实归因为 launch-bound（逐请求 Python 循环使 GPU 名义满载而带宽 util 仅 10--11\%），跨请求批量化已验证 3.8$\times$ 为修复方向。8B 稳态（bs16 $\times$ S=30K $\times$ n=256）：prefill 733.9 $\rightarrow$ 183.4s（4.00$\times$），decode step 65.0 $\rightarrow$ 33.5ms（1.94$\times$，反超 dense 的 40.4ms）。

\subsection{Profiling 方法论发现}

逐 kernel microbench 与生产 e2e 存在系统性口径差：select 热路径的 3 处 host 同步（\texttt{.item()} 与 \texttt{.any()}）在生产形态造成约 10 万次队列排空，合成 bench 无 GPU 队列时完全测不出，消除后同型负载 prefill 降低 6.7\%。同步消除类优化必须在 e2e 层验证。我们将此与离线重放排名不可替代 e2e 判决的观察合并为一条方法论：kernel 级与 e2e 两层评测都测、都诚实报告。

% =====================================================================
% 4 Related Work：三类聚类+定位句（MoBA 式 Category → Limitation → Position）。
% =====================================================================
\section{相关工作}

\textbf{稀疏注意力索引。}Quest 用 page 采样构造下界近似，KV 压缩类（SnapKV/H2O/PyramidKV）在 prefill 时一次性丢弃 KV。TLI 与它们的差异在三点：分数表示取上界（vs 下界——上界家族在 RULER 总 AVG 上领先 Quest $+8.40$，是长上下文质量分化的主要根源）、预算组织取分区（vs 全池）、两级级联（vs 单级）。我们在统一口径适配 SnapKV/H2O/PyramidKV 时发现：无显式 sink 保送的实现在 Qwen3-8B 上首 token 即 EOS——sink 在浅层不获选择信号但承载深层 0.59 的 attention mass，逐层恶性畸变。这既是对这类 baseline 适配的公平性修正，也旁证 sink 保送的必要性。

\textbf{训练类索引。}DSA~\cite{dsa}（DeepSeek V3.2）用 $I=\sum w\cdot\mathrm{ReLU}(q\cdot k)$ 训练选择器，64 头 $\times$128 维 FP8 top-2048，warm-up 1000 步；MoBA~\cite{moba} 用门控块稀疏。TLI training-free/weight-free，与其在 kernel 级统一 harness 对比。两者构成方法论分界：DSA 用训练换选择能力，可学任意投影；training-free 只能依赖与查询分布无关的位置先验——trace 离线成立而 e2e 崩坏的口径鸿沟正是这条分界线的实证刻画。

\textbf{两级结构。}HISA（COLM 2026）验证 block-coarse + token-refine 形态有效——「two-level」本身不是本文创新。TLI 的增量是两个正交维度的系统刻画：表示（上界 vs 下界）与预算组织（分区 vs 全池）的质量-速度影响。

\textbf{推理系统。}FlashAttention 系优化计算核，PagedAttention/vLLM 优化显存组织。TLI 与它们正交，在 sglang 生产形态集成。

% =====================================================================
% 5 Limitations：诚实边界。
% =====================================================================
\section{局限与未来工作}

精度侧：LongBench 首超 FullKV 的幅度在噪声边缘（+0.18）；RULER 最优臂 87.93 仍低于 FullKV 88.71（$-$0.78），残余代价集中于 multikey\_3（16K：75--80 vs TIA 89）——4bit 预算粒度在超多干扰 key 任务上主导。速度侧：多请求批量形态仍差 dense 8\%（launch-bound），跨请求批量化 forward\_extend 为明确修复方向；三方 indexer 延迟排名当前第三，跨请求批量化的 kernel 化为兑现路径。规模侧：子空间机制在 Qwen3-32B 上强度递减（entry recall 口径 0.66--0.87 vs 全维差距 $\leq$0.08，机制存在、如实报告）；MLA 架构的 latent 索引（DeepSeek 系）是 future work。gate 的 far 统计为 per-layer 单值，混合 batch 下跨请求污染（生产语义应 per-row）。

% =====================================================================
% 附录
% =====================================================================
\section*{附录 A：实验编号总索引}
全部实验编号与系统里程碑的脚本到结论映射（含 commit 号、对拍口径、负结果入册）见随源码发布的配套实验日志。

\section*{附录 B：可复现性}
H20-3e $\times$2（141GB），torch 2.8.0+cu128，sglang two-level-indexer 分支；RULER 官方预生成（KVCache-Factory 镜像）/ LongBench v1 全量 / Qwen3-8B、30B-A3B、32B 真实权重；string\_match\_all 官方口径打分；e2e 差分法采用稀疏/稠密配置交替运行的对照方式。

% =====================================================================
% References。作者列表不全的条目留 TODO，投稿前核对补全。
% =====================================================================
\begin{thebibliography}{99}\small

\bibitem{dsr1}
DeepSeek-AI. DeepSeek-R1: Incentivizing Reasoning Capability in LLMs via Reinforcement Learning. arXiv:2501.14348, 2025.

\bibitem{qwen3}
Qwen Team. Qwen3 Technical Report. arXiv:2505.09388, 2025.

\bibitem{fa}
T.~Dao, D.~Y. Fu, S.~Ermon, A.~Rudra, and C.~R\'{e}. FlashAttention: Fast and Memory-Efficient Exact Attention with IO-Awareness. In \emph{NeurIPS}, 2022.

\bibitem{vllm}
W.~Kwon, Z.~Li, S.~Zhuang, Y.~Sheng, L.~Zheng, C.~H. Yu, J.~E. Gonzalez, H.~Zhang, and I.~Stoica. Efficient Memory Management for Large Language Model Serving with PagedAttention. In \emph{SOSP}, 2023.

\bibitem{quest}
Y.~Tang, L.~Lin, Y.~Cao, S.~Li, F.~Yang, Y.~Wang, and T.~Wu. Quest: Query-Aware Sparsity for Efficient Long-Context LLM Inference. In \emph{ICLR}, 2025.

\bibitem{snapkv}
Y.~Li, Y.~Huang, B.~Yang, Y.~Li, Y.~Liu, and M.~Qiu. SnapKV: LLM Knows What You Are Looking For Before Generation. In \emph{NeurIPS}, 2024.

\bibitem{h2o}
Z.~Zhang, Y.~Sheng, T.~Zhou, D.~Chen, L.~Zheng, B.~Chen, L.~Cai, M.~Yan, Y.~Yu, J.~Tang, J.~Xu, P.~Zhang, and R.~Chen. H2O: Heavy-Hitter Oracle for Efficient Generative Inference of Large Language Models. In \emph{NeurIPS}, 2023.

\bibitem{pyramidkv}
Z.~Cai, Y.~Zhang, B.~Wang, X.~Liu, H.~Zhao, Y.~Lu, Y.~Wang, X.~Sun, K.~Liu, and Z.~Duan. PyramidKV: Dynamic KV Cache Compression Based on Pyramidal Information Funneling. arXiv:2406.02069, 2024.

\bibitem{streamingllm}
G.~Xiao, Y.~Tian, B.~Chen, S.~Han, and M.~Lewis. Efficient Streaming Language Models with Attention Sinks. In \emph{ICLR}, 2024.

\bibitem{longformer}
I.~Beltagy, M.~E. Peters, and A.~Cohan. Longformer: The Long-Document Transformer. arXiv:1904.07785, 2020.

\bibitem{bigbird}
M.~Zaheer, M.~Guruganesh, K.~A. Dubey, J.~Ainslie, C.~Alberti, S.~Onta\~{n}\'{o}n, P.~Pham, A.~Ravula, Q.~Wang, L.~Yang, and A.~Ahmed. Big Bird: Transformers for Longer Sequences. In \emph{NeurIPS}, 2020.

\bibitem{sparq}
% TODO: 投稿前核对作者列表
SparQ Attention: Accelerating LLM Inference with Accurate KV Retrieval, Minimized Decoding Costs, and No Fine-Tuning. arXiv:2406.09050, 2024.

\bibitem{dsa}
% TODO: 投稿前核对正式条目
DeepSeek-AI. DeepSeek-V3.2-Exp Technical Report: Boosting Long-Context Efficiency with Exact Sparse Attention. 2025.

\bibitem{moba}
% TODO: 投稿前核对作者列表
Moonshot AI. MoBA: Mixture of Block Attention for Efficient Long-Context LLMs. arXiv:2502.13189, 2025.

\bibitem{hisa}
% TODO: 投稿前核对作者列表
HISA: Efficient Hierarchical Indexing for Fine-Grained Sparse Attention. In \emph{COLM}, 2026. arXiv:2603.28458.

\bibitem{ruler}
C.-P. Hsieh, S.~Sun, S.~Kriman, S.~Ganti, J.~Zhang, B.~Kulis, K.~Papayiannis, M.~Fofanos, N.~Rudrajit, and C.~Tigges. RULER: What's the Real Context Size of Your Long-Context Language Models? arXiv:2404.06654, 2024.

\bibitem{longbench}
Y.~Bai, X.~Lv, J.~Zhang, H.~Zhao, X.~Qi, L.~Huang, Z.~Du, X.~Wang, Y.~Zheng, Y.~Zheng, J.~Zhang, Y.~Li, T.~Su, and J.~Tang. LongBench: A Bilingual, Multitask Benchmark for Long Context Understanding. In \emph{ACL}, 2024.

\end{thebibliography}

\end{document}
